
%(BEGIN_QUESTION)
% Copyright 2013, Tony R. Kuphaldt, released under the Creative Commons Attribution License (v 1.0)
% This means you may do almost anything with this work of mine, so long as you give me proper credit

In the mid 1800's, James Prescott Joule performed an experiment demonstrating how mechanical energy could be converted into heat.  His experimental apparatus is shown in this illustration, taken from issue 231 of {\it Harper's New Monthly Magazine} in August of 1869:

$$\epsfxsize=6in \includegraphics[width=15.5cm]{i02768x01.eps}$$

Examine this mechanism and then explain how it functions.  Specifically, what data would the experimenter record when running the experiment, and how would this data be used to establish a relationship between mechanical energy and heat?

\underbar{file i02768}
%(END_QUESTION)





%(BEGIN_ANSWER)

The falling weight pulls the string and turns the paddles inside the water-filled container, thereby causing the temperature of the water to rise.  Relevant data include:

\begin{itemize}
\item{} Initial temperature of the water
\item{} Final temperature of the water 
\item{} {\it Calculated water temperature rise} $\Delta T = T_f - T_o$
\vskip 5pt
\item{} Mass of water 
\item{} Specific heat of water 
\item{} Mass of paddle assembly
\item{} Specific heat of paddle assembly
\item{} Mass of thermometer
\item{} Specific heat of thermometer
\vskip 5pt
\item{} Magnitude of the weight
\item{} Initial height of the weight
\item{} Final height of the weight
\item{} {\it Calculated displacement} $\Delta x = x_f - x_o$
\end{itemize}

Relevant equations:

$$W = F \Delta x  \hskip 30pt  Q = mc \Delta T$$


%(END_ANSWER)





%(BEGIN_NOTES)


%INDEX% Physics, heat and temperature: equivalence of mechanical work with heat

%(END_NOTES)


